\section{Pillar 2 --- Kubernetes Performance Overhead at the Infrastructure Level}
\textit{(Theoretical basis for \sq{3})}

A substantial body of research has evaluated the performance overhead introduced by
containerisation and Kubernetes orchestration compared to bare-metal or \vm{} execution
\citep{cilic2023, liu2024}.

The primary sources of overhead in a Kubernetes execution model relevant to this study are
pod scheduling latency --- the time between a job being submitted and a pod being assigned
to a node --- and container startup time --- the time between pod assignment and the first
application instruction executing. These overheads are well-documented at the infrastructure
level but have not been measured at the application-level workflow execution layer for
pipeline orchestration systems specifically.

Additionally, Kubernetes introduces control plane communication overhead and network
namespace isolation costs that do not exist in a native \vm{} process. For short-duration
jobs, these fixed costs represent a proportionally larger fraction of total execution time
than for long-running jobs.

\textbf{Research connection:} \sq{3} extends infrastructure-level overhead measurement to
the application workflow execution layer and identifies at what concurrency level this
overhead is outweighed by the reliability gains from pod isolation --- the crossover point.
